\documentclass{article}

\usepackage{amsmath,amssymb}
\usepackage{xurl}
\usepackage[colorlinks=true,linkcolor=blue,citecolor=blue,urlcolor=blue]{hyperref}

% Shared commands for notes in docs/paper-gaps/.

\DeclareUnicodeCharacter{2080}{\ifmmode{}_{0}\else\textsubscript{0}\fi}
\DeclareUnicodeCharacter{2081}{\ifmmode{}_{1}\else\textsubscript{1}\fi}
\DeclareUnicodeCharacter{2082}{\ifmmode{}_{2}\else\textsubscript{2}\fi}
\DeclareUnicodeCharacter{2099}{\ifmmode{}_{n}\else\textsubscript{n}\fi}

\newcommand{\Lean}{\textsc{Lean}}
\ExplSyntaxOn
\NewDocumentCommand{\leanid}{m}{
  \ifmmode
    \textsf{\detokenize{#1}}
  \else
    \group_begin:
    \urlstyle{sf}
    \tl_set:Nn \l_tmpa_tl {#1}
    \tl_replace_all:Nnn \l_tmpa_tl {\_} {_}
    \exp_args:NV \path \l_tmpa_tl
    \group_end:
  \fi
}
\ExplSyntaxOff
\providecommand{\tr}{\operatorname{tr}}
\newcommand{\tnleanIssueBase}{https://github.com/LionSR/TNLean/issues/}
\newcommand{\tnleanPrBase}{https://github.com/LionSR/TNLean/pull/}
\newcommand{\tnleanDocBase}{https://sirui-lu.com/TNLean/docs/}
\newcommand{\ghissue}[1]{\href{\tnleanIssueBase#1}{GitHub issue \##1}}
\newcommand{\ghpr}[1]{\href{\tnleanPrBase#1}{PR \##1}}
\newcommand{\doclink}[1]{\href{\tnleanDocBase#1.html}{\path{#1}}}

% Dirac notation and the identity operator, as in blueprint/src/macros/common.tex.
% \providecommand keeps note-local definitions authoritative.
\usepackage{dsfont}
\providecommand{\ket}[1]{|#1\rangle}
\providecommand{\bra}[1]{\langle #1|}
\providecommand{\braket}[2]{\langle #1 | #2 \rangle}
\providecommand{\ketbra}[2]{|#1\rangle\!\langle #2|}
\providecommand{\Id}{\mathds{1}}
\providecommand{\one}{\mathds{1}}
\providecommand{\C}{\mathbb{C}}
\providecommand{\MN}[1]{M_{#1}(\C)}
\providecommand{\MD}{M_D(\C)}

% External referencing for paper sources.  The build helper writes sanitized
% auxiliary files under build/paper-aux/ and excludes duplicated source labels.
\usepackage{xr}

% Import a generated paper auxiliary file with a caller-supplied label prefix.
% The candidate paths support builds from docs/paper-gaps/, the repository root,
% and build/paper-aux/ itself.
\newcommand{\importpaperaux}[2]{%
  \IfFileExists{../../build/paper-aux/#2.aux}{%
    \externaldocument[#1]{../../build/paper-aux/#2}%
  }{%
    \IfFileExists{build/paper-aux/#2.aux}{%
      \externaldocument[#1]{build/paper-aux/#2}%
    }{%
      \IfFileExists{#2.aux}{%
        \externaldocument[#1]{#2}%
      }{}%
    }%
  }%
}

% General external reference macros
\newcommand{\paperref}[2]{\ref{#1-#2}}
\newcommand{\papereqref}[2]{(\ref{#1-#2})}

% CPSV16 source (1606.00608 / MPDO-22-12-17-2.tex) external document and shortcuts
\importpaperaux{cpsv16-}{1606.00608/MPDO-22-12-17-2-tnlean}
\newcommand{\cpsvref}[1]{\paperref{cpsv16}{#1}}
\newcommand{\cpsveqref}[1]{\papereqref{cpsv16}{#1}}

% Verdict marker: \gapnote{<kind>}{<status>}. Kind is the policy
% classification (clarification, local-correction, scope-restriction,
% unfaithful, false-source, open-gap); status is open, wip (elimination
% underway), resolved, or historical. Machine-read by the site index;
% prints a verdict line.
\newcommand{\gapnote}[2]{\par\noindent\textbf{Verdict:} #1 (#2).\par}


\title{Global and Per-Sector Unit-Weight Witnesses in the CPSV16 BNT Argument}
\author{}
\date{2026-05-30}

\begin{document}
\maketitle
\gapnote{scope-restriction}{resolved}

\section{Closure record}

Cirac--Perez-Garcia--Schuch--Verstraete 2016 normalizes the raw coefficients in
the basis-of-normal-tensors expansion by a single global condition: Section
II.C, line 246 requires \(|\mu_{j,q}|\le1\) for all copy weights and
\(|\mu_{j_\ast,q_\ast}|=1\) for some copy.  The paper-level canonical form does
not state that every BNT sector contains a copy of modulus one; a sector may
carry only strictly subunit copy weights while another sector holds the global
unit witness.  The formal predicate
\leanid{MPSTensor.IsBNTCanonicalForm} records exactly this global witness
through the field \leanid{MPSTensor.IsBNTCanonicalForm.weight_unit_exists}.

Earlier full-basis matching and global-gauge theorems assumed the stronger
per-sector condition \(\forall j,\ \exists q,\ |\mu_{j,q}|=1\), the price of an
asymptotic matching route that projected onto single sectors through the
Cesaro non-decay lemma.  That restriction has been eliminated.  The exact
linear-independence matcher
(\leanid{MPSTensor.exists_block_match_exact} for the equal case and
\leanid{MPSTensor.exists_block_match_exact_of_eventuallyProportional} for the
proportional case) works at one fixed sufficiently large length, and
\leanid{MPSTensor.IsBNTCanonicalForm.coeff_not_eventually_zero} shows that the
sector coefficient sum \(c^{(j)}_N=\sum_q\mu_{j,q}^N\) is nonzero infinitely
often for any nonzero copy weights, with no modulus assumption.  The matching
and global-gauge declarations therefore carry exactly the source hypothesis
set.  The unitary equal-case assembly of CPSV16 Corollary~C.5 is complete on
the same surface
(\leanid{MPSTensor.ft_sector_bnt_equal_unitary_global_gauge_witnesses} and
\leanid{MPSTensor.fundamentalTheorem_equal_canonicalForm_unitary}), again
without a per-sector unit-weight assumption.

The faithful structural normalization remains the global unit-weight field
\path{MPSTensor.IsBNTCanonicalForm.weight_unit_exists} in
\path{TNLean/MPS/FundamentalTheorem/SectorBNT/Basic.lean}.  The interaction
between this convention and the scale fixed by the MPDO renormalization
equations is the separate note
\path{docs/paper-gaps/cpsv16_unit_weight_rfp_scale_tension.tex}.

\bibliographystyle{alpha}
\bibliography{references}

\end{document}
